\documentclass[11pt]{article}

\usepackage[margin=1in]{geometry}
\usepackage{amsmath,amssymb,amsthm,mathtools}
\usepackage{enumitem}
\usepackage[colorlinks=true,linkcolor=blue,citecolor=blue,urlcolor=blue]{hyperref}

\newtheorem{theorem}{Theorem}
\newtheorem{corollary}[theorem]{Corollary}
\newtheorem{lemma}[theorem]{Lemma}
\theoremstyle{definition}
\newtheorem{definition}{Definition}[section]
\newtheorem*{fact}{Fact}
\newtheorem*{remark}{Remark}

\newcommand{\MSO}{\ensuremath{\mathrm{MSO}}}
\newcommand{\MSOtwo}{\ensuremath{\mathrm{MSO}_2}}
\newcommand{\defeq}{\coloneqq}
\newcommand{\Nodes}{N}
\newcommand{\Bags}{\mathrm{BAGS}}
\newcommand{\adh}{\mathrm{adh}}
\newcommand{\cone}{\mathrm{cone}}
\newcommand{\topnode}{\mathrm{top}}
\newcommand{\child}{\mathrm{child}}
\newcommand{\parent}{\mathrm{parent}}
\newcommand{\rootnode}{\mathrm{root}}
\newcommand{\legal}{\mathrm{legal}}
\newcommand{\vtx}{\mathrm{vtx}}
\newcommand{\set}{\mathrm{set}}
\newcommand{\adj}{\mathrm{adj}}
\newcommand{\Pred}{\mathcal P}
\newcommand{\Vrt}{\mathsf{Vert}}
\newcommand{\EObj}{\mathsf{EdgeObj}}
\newcommand{\eq}{\mathrm{eq}}
\newcommand{\cont}{\mathrm{cont}}
\newcommand{\emptysetf}{\mathrm{empty}}
\newcommand{\subsetf}{\mathrm{subset}}
\newcommand{\nonempty}{\mathrm{nonempty}}
\newcommand{\conn}{\mathrm{connected}}
\newcommand{\dang}{\mathrm{dangle}}
\newcommand{\Lab}{\mathrm{Lab}}
\newcommand{\Aut}{\mathcal{A}}
\newcommand{\calT}{\mathcal{T}}
\newcommand{\calC}{\mathcal{C}}
\newcommand{\calL}{\mathcal{L}}
\newcommand{\pow}{\mathcal{P}}
\newcommand{\tuple}[1]{\overline{#1}}
\newcommand{\st}{\;:\;}

\title{Expanded Notes on Courcelle's Theorem via Tree-Decomposition Encodings}
\author{Organized from \texttt{lecture\_note.pdf}; omitted proofs filled in}
\date{}

\begin{document}
\maketitle

\begin{abstract}
These notes reorganize the lecture note \texttt{lecture\_note.pdf} in proof order.
The main line is the standard route: encode a bounded-width, colored tree-decomposition as a finite-alphabet tree, interpret graph \MSO{} in tree \MSO{},
and use the equivalence between \MSO{}-definability and regular tree languages.
All lemmas that were stated as straightforward, omitted, or left as exercises in
the lecture note are expanded here.
\end{abstract}

\tableofcontents

\section{Courcelle's theorem}

Let graphs be finite simple graphs unless stated otherwise.  Fix a finite set
\(\Pred\) of unary predicate symbols and write
\[
  \tau_{\Pred}\defeq \{\adj\}\cup\Pred .
\]
A \(\tau_{\Pred}\)-graph is a finite simple graph together with an interpretation
\(\adj^G\subseteq V(G)^2\) of adjacency and an interpretation
\(P^G\subseteq V(G)\) for each \(P\in\Pred\).  Its treewidth is the treewidth of
the underlying graph, since unary predicates add no edges to the Gaifman graph.
The ordinary graph case is \(\Pred=\varnothing\).

Let us spell out the logical convention.  \MSO{} over \(\tau_{\Pred}\) is
one-sorted: first-order variables range over vertices, monadic second-order
variables range over sets of vertices, and the atomic formulas are
\(\adj(x,y)\), \(P(x)\) for \(P\in\Pred\), equality, and membership \(x\in X\).
By contrast, \MSOtwo{} on an original graph can quantify not only over vertices
and vertex sets, but also over edge-objects and sets of edge-objects.  It has an
incidence predicate \(\operatorname{inc}(v,e)\) connecting a vertex \(v\) to an
edge-object \(e\).  Thus an edge-object in \MSOtwo{} is not the same thing as
the binary adjacency relation \(\adj(x,y)\).  The \MSOtwo{} corollary below
therefore passes to a coloured incidence graph, where the unary predicates
\(\Vrt\) and \(\EObj\) remember which elements are original vertices and which
are original edge-objects.

\begin{theorem}[Courcelle's theorem for unary-expanded graphs]\label{thm:courcelle}
For every \(\omega\in\mathbb N\) and every fixed finite unary vocabulary
\(\Pred\), there exists an algorithm which, on input a \(\tau_{\Pred}\)-graph
\(G\), a binary tree-decomposition \((T,\beta)\) of the underlying graph of \(G\)
of width at most \(\omega\), and an \MSO{} sentence \(\varphi\) over
\(\tau_{\Pred}\), decides whether \(G\models\varphi\) in time
\[
  f_{\Pred}(\omega,|\varphi|)\cdot\bigl(|V(G)|+|N(T)|\bigr).
\]
In particular, for fixed \(\varphi\) and normalized decompositions with
\(|N(T)|=O(|V(G)|)\), the model-checking phase is linear in \(|V(G)|\).
\end{theorem}

The proof below constructs a finite alphabet \(\Sigma_\omega\) whose letters
remember both adjacency and unary predicate information, translates the input
into a \(\Sigma_\omega\)-tree, translates \(\varphi\) into a tree \MSO{} sentence,
and then runs the finite tree automaton corresponding to that tree sentence.

\begin{remark}
If one wants a statement with a supplied decomposition but with the shorter
linear bound \(f(\omega,\varphi)|V(G)|\), one should either assume that the input
decomposition has already been normalized so that \(|N(T)|=O(|V(G)|)\), or include
the linear-time normalization step in the algorithm.  The dependence on the
formula is written as \(|\varphi|\) because the algorithm must read and compile
the sentence; when \(\varphi\) is fixed this is absorbed into the parameter
function.
\end{remark}

\begin{theorem}[Computing treewidth]
There is an algorithm which, given a graph \(G\) and \(\omega\in\mathbb N\),
either outputs a tree-decomposition of width at most \(2\omega+1\), or correctly
decides that the treewidth of \(G\) is greater than \(\omega\), in time
\(2^{O(\omega)}|V(G)|\).
\end{theorem}

\begin{corollary}
For every fixed finite unary vocabulary \(\Pred\), there is an algorithm which,
given a \(\tau_{\Pred}\)-graph \(G\) whose underlying graph has treewidth at most
\(\omega\), and an \MSO{} sentence \(\varphi\) over \(\tau_{\Pred}\), decides
whether \(G\models\varphi\) in time \(f_{\Pred}(\omega,|\varphi|)|V(G)|\).
\end{corollary}

\begin{proof}
Use the treewidth approximation theorem to obtain a decomposition whose width is
bounded by a function of \(\omega\).  If desired, convert it to a binary or nice
binary tree-decomposition with only a linear blow-up.  Apply Theorem~\ref{thm:courcelle} with the
new width bound.  Since \(\omega\) is a parameter, changing the width bound only
changes the computable parameter function \(f\).
\end{proof}

\begin{corollary}[\MSOtwo{} properties]
Let \(\mathcal G\) be an \MSOtwo{}-definable graph property.  On every class of
graphs of treewidth at most \(\omega\), membership in \(\mathcal G\) can be
decided in time \(f(\omega,\mathcal G)|V(G)|\).
\end{corollary}

\begin{proof}
Let \(\psi\) be an \MSOtwo{} sentence defining \(\mathcal G\).  Given a graph
\(G\), form its \emph{coloured incidence graph} \(\widehat I(G)\).  Its universe
is the disjoint union \(V(G)\sqcup E(G)\).  The binary relation \(\adj\) is the
ordinary incidence adjacency: \(v\in V(G)\) is adjacent to \(e\in E(G)\) exactly
when \(v\) is an endpoint of \(e\).  In addition, \(\widehat I(G)\) carries two
unary predicates:
\[
  \Vrt(z) \quad\text{iff } z\in V(G),
  \qquad
  \EObj(z) \quad\text{iff } z\in E(G).
\]
Thus \(\widehat I(G)\) is a \(\tau_I\)-graph for
\(\tau_I=\{\adj,\Vrt,\EObj\}\).  The predicates \(\Vrt\) and \(\EObj\) are part
of the structure; they are not reconstructed from the uncoloured incidence
adjacency.

Translate \(\psi\) into a one-sorted \MSO{} sentence \(\psi^\dagger\) over
\(\tau_I\).  Atomic incidence becomes
\[
  \operatorname{inc}(x,e)^\dagger \defeq \adj(x,e),
\]
while typed quantifiers are guarded by the unary predicates:
\[
\begin{aligned}
  (\exists x\!:\!V\;\theta)^\dagger
    &\defeq \exists x\bigl(\Vrt(x)\wedge\theta^\dagger\bigr),\\
  (\exists e\!:\!E\;\theta)^\dagger
    &\defeq \exists e\bigl(\EObj(e)\wedge\theta^\dagger\bigr).
\end{aligned}
\]
For set variables, use the guards
\[
  \operatorname{VertSet}(X)\defeq
    \forall z(z\in X\to\Vrt(z)),\qquad
  \operatorname{EdgeSet}(F)\defeq
    \forall z(z\in F\to\EObj(z)),
\]
and translate vertex-set and edge-set quantifiers by adding these guards.
Universal quantifiers are translated dually.  A structural induction on
\(\psi\) gives
\[
  G\models_{\MSOtwo}\psi
  \quad\Longleftrightarrow\quad
  \widehat I(G)\models_{\MSO}\psi^\dagger .
\]

It remains to check bounded treewidth and size.  If \((T,\beta)\) has width
\(\omega\), then the underlying incidence graph of \(\widehat I(G)\) has width at
most \(\max(\omega,2)\): for every edge \(e=uv\), choose a bag containing \(u,v\)
and attach a new leaf bag \(\{u,v,e\}\).  The old bags cover the original
vertices; the new leaf bags cover incidence edges; the bags containing every
old vertex remain connected; and each edge-object appears in one new leaf bag.
Unary predicates do not affect treewidth.  Also, a simple graph of treewidth at
most \(\omega\) is \(\omega\)-degenerate, hence has \(O(\omega|V(G)|)\) edges.
Therefore \(|V(\widehat I(G))|=|V(G)|+|E(G)|=O_\omega(|V(G)|)\).

Apply the previous corollary to the \(\tau_I\)-graph \(\widehat I(G)\) and the
sentence \(\psi^\dagger\).  Since \(\psi\) is fixed by the property
\(\mathcal G\), the dependence on \(\psi^\dagger\) is absorbed into
\(f(\omega,\mathcal G)\).
\end{proof}

\section{Encoding a colored decomposition as a \texorpdfstring{\(\Sigma\)}{Sigma}-tree}

For the rest of the proof, fix the finite unary vocabulary \(\Pred\).  Fix a
rooted binary tree-decomposition \((T,\beta)\) of the underlying graph of a
\(\tau_{\Pred}\)-graph \(G\), of width at most \(\omega\).  The set of nodes of
\(T\) is denoted by \(N(T)\).  If \(t\) is not the root, \(p(t)\) denotes its
parent.  Define
\[
  \adh(t) \defeq
  \begin{cases}
    \emptyset, & t \text{ is the root},\\
    \beta(t)\cap\beta(p(t)), & \text{otherwise},
  \end{cases}
\]
and let \(\cone(t)\) be the union of the bags in the subtree rooted at \(t\).

For a vertex \(v\in V(G)\), let
\[
  \Bags(v)\defeq \{t\in N(T): v\in \beta(t)\}.
\]
Because \((T,\beta)\) is a tree-decomposition, \(\Bags(v)\) is a nonempty
connected set of tree nodes.  For a nonempty connected node set \(U\), let
\(\topnode(U)\) be the unique node of \(U\) closest to the root; since \(U\) is
connected, every node of \(U\) is a descendant of \(\topnode(U)\).

\begin{fact}[Bag-injective coloring]
There is a coloring \(c:V(G)\to C_\omega\defeq\{0,\ldots,\omega\}\) such that no
bag contains two vertices of the same color.
\end{fact}

\begin{proof}
Abbreviate \(\tau(v)\defeq\topnode(\Bags(v))\), the topmost node of the bag set of
\(v\).  Recall that \(\Bags(v)\) is nonempty and connected, so \(\tau(v)\) is
well-defined and every node of \(\Bags(v)\) is a descendant of \(\tau(v)\).

Color the vertices one at a time, in any order of nondecreasing depth of
\(\tau(v)\) (equivalently, sweep the nodes of \(T\) from the root downward and
color \(v\) when \(\tau(v)\) is reached).  When \(v\) is colored at
\(t\defeq\tau(v)\), assign it any color in \(C_\omega\) not already used on the
colored vertices of \(\beta(t)\).

\emph{Claim.} At the moment \(v\) is colored, every already-colored vertex \(u\)
with \(\Bags(u)\cap\Bags(v)\ne\emptyset\) lies in \(\beta(t)\).

Pick \(s\in\Bags(u)\cap\Bags(v)\).  As \(s\in\Bags(v)\), the topmost node
\(t=\tau(v)\) is an ancestor of \(s\); as \(s\in\Bags(u)\), the node \(\tau(u)\)
is also an ancestor of \(s\).  Hence \(\tau(u)\) and \(t\) both lie on the path
from the root to \(s\), on which ancestry is a total order.  Since \(u\) was
colored before \(v\), we have \(\mathrm{depth}(\tau(u))\le\mathrm{depth}(t)\), so
\(\tau(u)\) is an ancestor of \(t\) (or equals it).  Then \(t\) lies on the path
from \(\tau(u)\) to \(s\), whose endpoints both lie in the connected set
\(\Bags(u)\); therefore \(t\in\Bags(u)\), i.e.\ \(u\in\beta(t)\).  This proves the
claim.

By the claim, the colors forbidden to \(v\) are those of the already-colored
vertices in \(\beta(t)\setminus\{v\}\), of which there are at most \(\omega\)
because \(|\beta(t)|\le\omega+1\).  Hence one of the \(\omega+1\) colors of
\(C_\omega\) is available, and the procedure never gets stuck.

The resulting coloring is bag-injective.  If \(u\ne v\) share a bag, then
\(\Bags(u)\cap\Bags(v)\ne\emptyset\), so whichever of the two is colored later finds the
other in its bag \(\beta(\tau(\cdot))\) by the claim and avoids its color; thus
\(c(u)\ne c(v)\).  The procedure colors each vertex once during a single
root-to-leaf sweep, so it is constructive.
\end{proof}

\subsection{Rooted unary-expanded graphs and gluing}

\begin{definition}
A \(k\)-rooted \(\tau_{\Pred}\)-graph is a triple \((H,R,\iota)\) where \(H\) is a
\(\tau_{\Pred}\)-graph, \(R\subseteq V(H)\) is the root or boundary, and
\(\iota\) is a partial injective map from \(V(H)\) to \([k]\) whose domain
contains \(R\).  When \(\iota(v)=i\), we say that \(v\) has label \(i\).
\end{definition}

Given \(k\)-rooted \(\tau_{\Pred}\)-graphs \((H_1,R_1,\iota_1)\) and
\((H_2,R_2,\iota_2)\), the partial gluing
\[
  (H_1,R_1,\iota_1)\oplus_k(H_2,R_2,\iota_2)
\]
is defined when the labelled rooted \(\tau_{\Pred}\)-graph induced by \(R_2\) is
label-preserving isomorphic to the induced labelled \(\tau_{\Pred}\)-subgraph of
\(H_1\) on vertices with the same labels.  Thus, on the common labelled root,
both the adjacency relation and every unary predicate \(P\in\Pred\) agree.  Then
we take the disjoint union and identify every \(v\in R_2\) with the unique vertex
\(v'\in V(H_1)\) satisfying \(\iota_1(v')=\iota_2(v)\).  The root and labelling of
the result are inherited from the first graph, and each unary predicate is
inherited from the two pieces.  The compatibility condition makes the predicate
interpretations well-defined on identified vertices.

For a node \(t\), define the \emph{node graph}
\[
  (H_t,R_t,\iota_t)
  \defeq (G[\beta(t)],\adh(t),c|_{\beta(t)}).
\]
The \emph{cone graph} at \(t\) is
\[
  (G_t,R_t,\iota_t) \defeq (G[\cone(t)],\adh(t),c|_{\beta(t)}),
\]
whose external root is the adhesion \(\adh(t)\) but which retains the bag labels
\(c|_{\beta(t)}\).  These temporary labels are exactly the ones supplied by the
node graph during gluing; only the labels on \(\adh(t)\) are used when the cone
is later attached to its parent.

We first record three consequences of the running-intersection axiom; they are
used in the cone-gluing lemma below and again in the decoding lemmas.

\begin{lemma}[Cone structure]\label{lem:cone-struct}
\begin{enumerate}[label=\textup{(\roman*)}]
  \item For every node \(t\) and every child \(s\) of \(t\),
    \(\cone(s)\cap\beta(t)=\adh(s)\).
  \item For distinct children \(t_1,t_2\) of a node \(t\),
    \(\cone(t_1)\cap\cone(t_2)\subseteq\adh(t_1)\cap\adh(t_2)\).
  \item Every edge of \(G[\cone(t)]\) has both endpoints in a common bag
    \(\beta(s)\) with \(s\) in the subtree rooted at \(t\).  When \(t\) is internal
    with children \(t_1,t_2\), it is therefore an edge of one of \(G[\beta(t)]\),
    \(G[\cone(t_1)]\), \(G[\cone(t_2)]\).
\end{enumerate}
\end{lemma}

\begin{proof}
(i) Since \(\adh(s)=\beta(s)\cap\beta(t)\) and \(\beta(s)\subseteq\cone(s)\), we
have \(\adh(s)\subseteq\cone(s)\cap\beta(t)\).  Conversely, let
\(v\in\cone(s)\cap\beta(t)\).  Then \(v\in\beta(s')\) for some node \(s'\) in the
subtree rooted at \(s\), so \(s'\in\Bags(v)\), and also \(t\in\Bags(v)\).  Because
\(s'\) is a descendant of \(s\) and \(t\) is the parent of \(s\), the path in
\(T\) from \(s'\) to \(t\) passes through \(s\).  As \(\Bags(v)\) is connected and
contains \(s'\) and \(t\), it contains this whole path, in particular \(s\); hence
\(v\in\beta(s)\), and \(v\in\beta(s)\cap\beta(t)=\adh(s)\).

(ii) Let \(v\in\cone(t_1)\cap\cone(t_2)\), say \(v\in\beta(s_1)\) with \(s_1\) in
the subtree rooted at \(t_1\) and \(v\in\beta(s_2)\) with \(s_2\) in the subtree
rooted at \(t_2\); thus \(s_1,s_2\in\Bags(v)\).  As \(t_1\ne t_2\) are distinct
children of \(t\), the subtrees rooted at \(t_1\) and \(t_2\) are disjoint, and the
unique path in \(T\) from \(s_1\) to \(s_2\) runs
\(s_1\to\cdots\to t_1\to t\to t_2\to\cdots\to s_2\), passing through \(t\).  Since
\(\Bags(v)\) is connected and contains \(s_1,s_2\), it contains \(t\), so
\(v\in\beta(t)\).  Thus \(\cone(t_1)\cap\cone(t_2)\subseteq\beta(t)\); intersecting
with part~(i) gives
\(\cone(t_1)\cap\cone(t_2)\subseteq\cone(t_i)\cap\beta(t)=\adh(t_i)\) for
\(i=1,2\).

(iii) Let \(uv\) be an edge of \(G[\cone(t)]\), so \(u,v\in\cone(t)\) and
\(uv\in E(G)\).  By the edge-covering axiom some bag \(\beta(s_0)\) contains both
\(u\) and \(v\).  If \(s_0\) lies in the subtree rooted at \(t\), take \(s=s_0\).
Otherwise \(s_0\) is not a descendant of \(t\); since \(u\in\cone(t)\), there is a
node \(a\in\Bags(u)\) in the subtree rooted at \(t\), and the path from \(a\) to
\(s_0\) leaves that subtree, hence passes through its root \(t\).  As \(\Bags(u)\)
is connected and contains \(a\) and \(s_0\), it contains \(t\), so \(u\in\beta(t)\);
the same argument gives \(v\in\beta(t)\), and \(s=t\) works.  In either case both
endpoints lie in \(\beta(s)\) with \(s\) in the subtree rooted at \(t\).  When
\(t\) is internal with children \(t_1,t_2\), the node \(s\) is \(t\), or lies in
the subtree rooted at \(t_1\), or in the subtree rooted at \(t_2\); correspondingly
\(\beta(s)\subseteq\beta(t)\), \(\beta(s)\subseteq\cone(t_1)\), or
\(\beta(s)\subseteq\cone(t_2)\), and the edge belongs to the induced subgraph on
that vertex set.
\end{proof}

\begin{lemma}[Cone gluing]\label{lem:cone-gluing}
Let \(t\) be an internal node with children \(t_1,t_2\).  If
\((G_i,R_i,\iota_i)\) is the cone graph of \(t_i\), then
\[
  (G_t,R_t,\iota_t)
  = (H_t,R_t,\iota_t)\oplus (G_1,R_1,\iota_1)\oplus (G_2,R_2,\iota_2).
\]
\end{lemma}

\begin{proof}
For each child \(t_i\), the root \(R_i=\adh(t_i)=\beta(t_i)\cap\beta(t)\) is a
subset of \(\beta(t)\).  Because the coloring is injective on \(\beta(t)\), every
label appearing in \(R_i\) identifies a unique vertex of \(H_t\).  The induced
substructure on \(R_i\) inside the child cone is the same as the induced
substructure on that same vertex set inside \(H_t\), since both are induced
\(\tau_{\Pred}\)-subgraphs of the original structure \(G\).  Thus the edge
relation and all unary predicates agree on the common labelled root.  Therefore
the first gluing is defined, and the same argument applies to the second child.

The vertex set of the result is
\(\beta(t)\cup\cone(t_1)\cup\cone(t_2)=\cone(t)\).  The gluing identifies a vertex
of a child cone with a vertex of \(H_t\) only when they share a label, that is,
along \(R_i=\adh(t_i)\).  By Lemma~\ref{lem:cone-struct}(i),
\(\cone(t_i)\cap\beta(t)=\adh(t_i)\), so these are exactly the vertices that
\(\cone(t_i)\) and \(\beta(t)\) have in common; and by Lemma~\ref{lem:cone-struct}(ii),
\(\cone(t_1)\cap\cone(t_2)\subseteq\adh(t_1)\cap\adh(t_2)\), so a vertex lying in
both child cones is, through its label, identified with the same vertex of
\(H_t\).  Hence the quotient identifies precisely the shared vertices of the three
pieces and nothing more.

Gluing creates no new edges; it takes the union of the edges already present in
the three induced graphs.  By Lemma~\ref{lem:cone-struct}(iii), every edge of
\(G[\cone(t)]\) is an edge of \(G[\beta(t)]\), \(G[\cone(t_1)]\), or
\(G[\cone(t_2)]\), hence of one of the pieces; conversely every edge of a piece is
an induced edge of \(G[\cone(t)]\).  The same union argument applies to each unary
predicate, since each piece carries the restriction of the predicate
interpretation from the same ambient structure \(G\).  The root of the result is
\(R_t=\adh(t)\) and, since gluing inherits the labelling of its first operand, its
labels are \(c|_{\beta(t)}\), matching the cone graph at \(t\).  Thus the glued
\(\tau_{\Pred}\)-graph equals \(G_t\), as claimed, up to the canonical
label-preserving isomorphism of the disjoint-union-and-quotient model of
\(\oplus\).
\end{proof}

\subsection{The alphabet and legal trees}

Let
\[
  \Sigma_\omega \defeq
  \{(H,R): H \text{ is a } \tau_{\Pred}\text{-graph with }
       V(H)\subseteq C_\omega,\ R\subseteq V(H)\}.
\]
A letter is a rooted \(\tau_{\Pred}\)-graph on a subset of the fixed color set
\(C_\omega\); the identity map is the implicit labelling.  Equivalently, a
letter records the induced adjacency relation on the bag and, for every present
color \(i\), which unary predicates \(P\in\Pred\) hold of the element of color
\(i\).  Since both \(C_\omega\) and \(\Pred\) are finite, the set
\(\Sigma_\omega\) is finite.

The encoding \((T,\rho)\) of \((G,c,(T,\beta))\) labels each node \(t\) by the
rooted \(\tau_{\Pred}\)-graph obtained from \((H_t,R_t,\iota_t)\) after renaming
every vertex \(v\in\beta(t)\) to its color \(c(v)\).  Since colors are injective
on bags, this renaming is well-defined and loses no information inside a single
bag.

\begin{definition}
A \(\Sigma_\omega\)-tree \((T,\rho)\) is \emph{legal} if its root \(r\) carries a
letter \(\rho(r)=(H_r,R_r)\) with \(R_r=\varnothing\), and if for every node
\(x\) with parent \(y\), writing \(\rho(x)=(H_x,R_x)\) and \(\rho(y)=(H_y,R_y)\),
we have
\[
  R_x\subseteq V(H_y)
  \quad\text{and}\quad
  H_x[R_x]=H_y[R_x].
\]
The equality here is equality of induced \(\tau_{\Pred}\)-graphs: the edge
relation and every unary predicate agree on the shared colors.
\end{definition}

\begin{lemma}
There is an \MSO{} sentence \(\varphi_{\legal}\) over the tree vocabulary
\(\tau_\Sigma=\{\child_1,\child_2\}\cup\{P_a:a\in\Sigma_\omega\}\) such that
\((T,\rho)\models\varphi_{\legal}\) iff \((T,\rho)\) is legal.
\end{lemma}

\begin{proof}
For letters \(a=(H_a,R_a)\) and \(b=(H_b,R_b)\), write
\(\operatorname{Comp}(a,b)\) when a child labelled \(a\) is compatible with a
parent labelled \(b\), namely
\[
  R_a\subseteq V(H_b)\quad\text{and}\quad H_a[R_a]=H_b[R_a].
\]
Again the equality is equality of induced \(\tau_{\Pred}\)-graphs, so unary
predicate tags must agree on every shared root color.  This is a finite relation
on \(\Sigma_\omega\).  Define
\[
  \parent(y,x) \defeq \child_1(y,x)\vee \child_2(y,x)
\]
and
\[
  \varphi_{\operatorname{loc}}(x,y)
  \defeq
  \parent(y,x)\wedge
  \bigvee_{\operatorname{Comp}(a,b)}(P_a(x)\wedge P_b(y)).
\]
Writing \(\Sigma_\omega^{\varnothing}\defeq\{(H,R)\in\Sigma_\omega:R=\varnothing\}\)
for the letters with empty root, set
\[
  \varphi_{\legal}
  \defeq
  \forall x\,\Bigl(\rootnode(x)\to\bigvee_{a\in\Sigma_\omega^{\varnothing}}P_a(x)\Bigr)
  \wedge
  \forall x\,\bigl(\rootnode(x)\vee \exists y\,\varphi_{\operatorname{loc}}(x,y)\bigr),
\]
where \(\rootnode(x)\) abbreviates \(\neg\exists y\,\parent(y,x)\).  Since every
node carries exactly one letter, the first conjunct says the root letter has
empty root set, and the finite disjunction in the second says exactly that the
letter of a non-root node is compatible with the letter of its parent.  This is
precisely legality.
\end{proof}

\subsection{Decoding a legal tree}

Suppose \((T,\rho)\) is legal, and write \(\rho(t)=(H_t,R_t)\) for the letter at
\(t\), so that \(V(H_t)\subseteq C_\omega\) and \(R_t\subseteq V(H_t)\).  For each
node \(t\) and color \(i\in V(H_t)\) introduce a formal \emph{slot} \(v(t,i)\).

The \emph{decoded graph} \(G\) glues these slots along the adhesions recorded by
the root sets.  On the set of all slots define the equivalence relation \(\sim\)
generated by
\[
  v(x,i)\sim v(p(x),i)\qquad\text{whenever \(x\) has parent \(p(x)\) and
  \(i\in R_x\).}
\]
Legality gives \(R_x\subseteq V(H_{p(x)})\), so \(v(p(x),i)\) is a slot whenever
\(v(x,i)\) is and \(i\in R_x\); each generator thus relates two genuine slots.  A
\emph{decoded vertex} is a \(\sim\)-class \([v(t,i)]\), and \(V(G)\) is the set of
classes.  Every generator preserves the color coordinate \(i\), so all slots of a
class share one color, and we set \(c([v(t,i)])\defeq i\).  Put
\[
  \beta(t)\defeq\{[v(t,i)]:i\in V(H_t)\},
\]
so that \(v\in\beta(t)\) iff some slot \(v(t,i)\) belongs to \(v\).  Adjacency and
the unary predicates are induced from the letters: \([v(t,i)]\) and \([v(s,j)]\)
are adjacent iff some node \(u\) carries slots \(v(u,i),v(u,j)\) of the two classes
with \(ij\in E(H_u)\), and \(P\) holds of \([v(t,i)]\) iff \(i\in P^{H_t}\) for a
representative.  Operationally, \(G\) is the bottom-up gluing of the letters by the
same operation \(\oplus\) used in the cone gluing of Lemma~\ref{lem:cone-gluing};
since the root letter has \(R_r=\varnothing\), no labels survive at the root and
\(G\) is an ordinary \(\tau_{\Pred}\)-graph.

For a color \(i\), let \(\Gamma_i\) be the subgraph of \(T\) whose vertices are the
nodes \(s\) with \(i\in V(H_s)\) and whose edges are the tree edges \(\{x,p(x)\}\)
with \(i\in R_x\).

\begin{lemma}\label{lem:decode-occ}
Let \((T,\rho)\) be legal and \(i\in C_\omega\).  Two slots \(v(s,i)\) and
\(v(t,i)\) are \(\sim\)-equivalent iff every edge of the unique \(s\)--\(t\) path in
\(T\) lies in \(\Gamma_i\), that is, carries color \(i\) in its child root.
Consequently each decoded vertex \(v\) of color \(i\) corresponds to a connected
component of \(\Gamma_i\), and its occurrence set
\(U_v\defeq\{s\in N(T):v\in\beta(s)\}\) is a nonempty connected subtree of \(T\).
\end{lemma}

\begin{proof}
Each generator of \(\sim\) relates two color-\(i\) slots across an edge
\(\{x,p(x)\}\) of \(\Gamma_i\), and conversely every edge of \(\Gamma_i\) is such a
generator.  Hence on color-\(i\) slots \(\sim\) is exactly the connectivity
relation of \(\Gamma_i\): \(v(s,i)\sim v(t,i)\) iff \(s\) and \(t\) lie in one
component of \(\Gamma_i\).  Because the path between two nodes of a tree is unique,
\(s\) and \(t\) are \(\Gamma_i\)-connected iff every edge of that path is a
\(\Gamma_i\)-edge, which is the stated path condition.

The color-\(i\) slots of a decoded vertex \(v\) are the vertices of one component
of \(\Gamma_i\), so \(U_v\) is that component's node set.  Every edge of
\(\Gamma_i\) is a tree edge, so the component is a connected subgraph of \(T\) built
from tree edges, hence a subtree; thus \(U_v\) is nonempty and connected.
\end{proof}

\begin{lemma}\label{lem:decode-wd}
If \((T,\rho)\) is legal, its decoding \((G,c,(T,\beta))\) is well-defined, and
\((T,\beta)\) is a tree-decomposition of the underlying graph of \(G\) on which
\(c\) is bag-injective.
\end{lemma}

\begin{proof}
The color \(c([v(t,i)])=i\) is well-defined because every generator preserves the
color coordinate.  The predicates are well-defined because a generator
\(v(x,i)\sim v(p(x),i)\) has \(i\in R_x\) and legality gives
\(H_x[R_x]=H_{p(x)}[R_x]\) as induced \(\tau_{\Pred}\)-graphs, so the color-\(i\)
predicate tags agree at \(x\) and \(p(x)\); this propagates along any chain of
generators.  The edge relation is well-defined as an existential over
representatives.

Each bag \(\beta(t)\) has at most \(\omega+1\) vertices since
\(V(H_t)\subseteq C_\omega\), and \(c\) is injective on \(\beta(t)\) because a
letter has at most one vertex of each color.  Every decoded vertex \([v(t,i)]\)
lies in \(\beta(t)\), so the bags cover \(V(G)\); and each edge of \(G\), arising
from some \(ij\in E(H_u)\), has both endpoints \([v(u,i)],[v(u,j)]\) in
\(\beta(u)\), so the bags cover every edge.  By Lemma~\ref{lem:decode-occ} the
occurrence set \(U_v\) of every decoded vertex is connected, the
running-intersection axiom.  Hence \((T,\beta)\) is a tree-decomposition, and \(c\)
is bag-injective by construction.
\end{proof}

\begin{lemma}\label{lem:decode-adh}
Let \((T,\rho)\) be legal and let \(x\) be a non-root node with parent
\(y=p(x)\).  Then, in the decoding,
\[
  \beta(x)\cap\beta(y)=\{[v(x,i)]:i\in R_x\}.
\]
\end{lemma}

\begin{proof}
If \(i\in R_x\), then \(i\in V(H_x)\) and, by legality, \(i\in V(H_y)\); the
generator \(v(x,i)\sim v(y,i)\) gives \([v(x,i)]=[v(y,i)]\), which lies in both
\(\beta(x)\) and \(\beta(y)\).  This proves \(\supseteq\).

Conversely, let \(w\in\beta(x)\cap\beta(y)\), say \(w=[v(x,i)]=[v(y,j)]\) with
\(i\in V(H_x)\) and \(j\in V(H_y)\).  Comparing colors gives \(i=j\), so
\(v(x,i)\sim v(y,i)\).  By Lemma~\ref{lem:decode-occ} the \(x\)--\(y\) path lies in
\(\Gamma_i\); this path is the single edge \(\{x,y\}\), which is a
\(\Gamma_i\)-edge exactly when \(i\in R_x\).  Hence \(i\in R_x\) and
\(w=[v(x,i)]\) with \(i\in R_x\).
\end{proof}

\begin{lemma}\label{lem:decode-bag}
Let \((T,\rho)\) be legal.  For every node \(t\), the map \(i\mapsto[v(t,i)]\) is
an isomorphism of \(\tau_{\Pred}\)-graphs from the letter \(H_t\) onto the induced
subgraph \(G[\beta(t)]\) of the decoded graph.
\end{lemma}

\begin{proof}
The map sends \(V(H_t)\) onto \(\beta(t)\) and is injective, since
\(c([v(t,i)])=i\) distinguishes the images; so it is a bijection
\(V(H_t)\to\beta(t)\).  It preserves the unary predicates by the definition of the
decoded predicates, \(P([v(t,i)])\) iff \(i\in P^{H_t}\).

For adjacency, fix \(i,j\in V(H_t)\).  If \(ij\in E(H_t)\), then taking \(u=t\) in
the definition of decoded adjacency shows \([v(t,i)]\) and \([v(t,j)]\) are
adjacent.  Conversely, suppose they are adjacent.  By definition there is a node
\(u\) with \(v(u,i)\sim v(t,i)\), \(v(u,j)\sim v(t,j)\), and \(ij\in E(H_u)\).  By
Lemma~\ref{lem:decode-occ} the \(u\)--\(t\) path lies in both \(\Gamma_i\) and
\(\Gamma_j\), so every edge \(\{z,p(z)\}\) on it has \(i,j\in R_z\).  Legality
gives \(H_z[R_z]=H_{p(z)}[R_z]\) as induced \(\tau_{\Pred}\)-graphs, whence
\(ij\in E(H_z)\) iff \(ij\in E(H_{p(z)})\) for these \(i,j\in R_z\).  Thus the
truth value of ``\(ij\in E(H_\cdot)\)'' is constant along the \(u\)--\(t\) path,
and \(ij\in E(H_u)\) yields \(ij\in E(H_t)\).  Hence the map preserves and reflects
adjacency, so it is an isomorphism onto \(G[\beta(t)]\).
\end{proof}

\begin{lemma}\label{lem:legal-iff-enc}
A \(\Sigma_\omega\)-tree \((T,\rho)\) is legal iff it is the encoding of some
\((\omega+1)\)-colored \(\tau_{\Pred}\)-graph with a binary tree-decomposition
of width at most \(\omega\).  Moreover, in the forward direction the witnessing
triple is the decoding of \((T,\rho)\).
\end{lemma}

\begin{proof}
If \((T,\rho)\) is an encoding, legality follows directly from the definition of
adhesion: the child root is exactly the intersection with the parent bag, and
both induced \(\tau_{\Pred}\)-graphs are induced subgraphs of the same structure
\(G\) on that intersection.

Conversely, assume \((T,\rho)\) is legal and decode it as \((G,c,(T,\beta))\).
By Lemma~\ref{lem:decode-wd}, \((T,\beta)\) is a tree-decomposition on which \(c\)
is bag-injective; each bag \(\beta(t)\) has at most \(\omega+1\) vertices because
\(V(H_t)\subseteq C_\omega\), so the width is at most \(\omega\), and \(T\) is
binary.  It remains to check that re-encoding recovers \((T,\rho)\).  Re-encoding
labels each node \(t\) by the letter obtained from \(G[\beta(t)]\), with root
\(\adh(t)\), by renaming each vertex to its color.  By Lemma~\ref{lem:decode-bag}
the renamed graph is exactly \(H_t\), with the same adjacency and unary
predicates.  For the root set: at the tree root \(\adh(t)=\varnothing\) and
legality forces \(R_t=\varnothing\); at a non-root node
\(\adh(t)=\beta(t)\cap\beta(p(t))\), which Lemma~\ref{lem:decode-adh} identifies
with \(\{[v(t,i)]:i\in R_t\}\), whose color set is \(R_t\).  Hence the re-encoded
letter at \(t\) is \((H_t,R_t)=\rho(t)\), so re-encoding the decoded triple returns
\((T,\rho)\).
\end{proof}

\section{Combinatorial properties of defining tuples}

Recall the bag set \(\Bags(v)=\{t\in N(T):v\in\beta(t)\}\), a nonempty connected
set of tree nodes, and the topmost node \(\topnode(U)\) of a nonempty connected
node set \(U\).  The \emph{defining pair} of \(v\) is \((\Bags(v),c(v))\).  A node
\(x\notin U\) is a \emph{dangle} of \(U\) if its parent lies in \(U\).

\begin{lemma}[Defining pairs]\label{lem:defining-pairs}
Let \((T,\beta)\) be a binary tree-decomposition of \(G\) of width at most
\(\omega\), and let \(c:V(G)\to C_\omega\) be injective on every bag.  For
\(U\subseteq N(T)\) and \(i\in C_\omega\), the pair \((U,i)\) is the defining pair
of some vertex of \(G\) iff the following conditions hold.
\begin{enumerate}[label=(\roman*)]
  \item \(U\) is nonempty and connected in \(T\).
  \item For every \(t\in U\), the bag \(\beta(t)\) contains a vertex of color \(i\).
  \item For every \(t\in U\setminus\{\topnode(U)\}\), the adhesion \(\adh(t)\)
    contains a vertex of color \(i\).
  \item The adhesion \(\adh(\topnode(U))\) contains no vertex of color \(i\).
  \item If \(x\notin U\) and the parent of \(x\) lies in \(U\), then \(\adh(x)\)
    contains no vertex of color \(i\).
\end{enumerate}
\end{lemma}

\begin{proof}
First suppose \((U,i)=(\Bags(v),c(v))\).  Then \(U\) is nonempty and connected by
the running-intersection condition of the tree-decomposition, and every bag in
\(U\) contains the color-\(i\) vertex \(v\).  If
\(t\in U\setminus\{\topnode(U)\}\), then the parent of \(t\) also lies on the
connected set \(U\), so \(v\in \beta(t)\cap\beta(p(t))=\adh(t)\).  If
\(v\in\adh(\topnode(U))\), then the parent of \(\topnode(U)\) would also contain
\(v\), contradicting the choice of the topmost node.  Finally, if a dangle
\(x\notin U\) had a color-\(i\) vertex in \(\adh(x)\), that vertex and \(v\) would
both lie in \(\beta(p(x))\) and have the same color.  Bag-injectivity gives that
this vertex is \(v\), whence \(x\in\Bags(v)=U\), a contradiction.

Conversely, assume (i)--(v).  For each \(t\in U\), let \(v_t\) be the unique
color-\(i\) vertex in \(\beta(t)\).  Uniqueness follows from bag-injectivity.  If
\(t\in U\setminus\{\topnode(U)\}\), condition (iii) gives a color-\(i\) vertex in
\(\beta(t)\cap\beta(p(t))\), so \(v_t=v_{p(t)}\).  Since \(U\) is connected, all
vertices \(v_t\) are a single vertex \(v\).  Thus \(U\subseteq\Bags(v)\).

If \(\Bags(v)\) had a node outside \(U\), the connected set \(\Bags(v)\) would
cross the boundary of \(U\).  It can cross upward only through the parent of
\(\topnode(U)\), which would put \(v\) in \(\adh(\topnode(U))\), contradicting
(iv).  It can cross downward only through a dangle \(x\notin U\) whose parent is
in \(U\), which would put \(v\) in \(\adh(x)\), contradicting (v).  Hence
\(\Bags(v)=U\) and \(c(v)=i\).
\end{proof}

\begin{lemma}[Distinct defining pairs]\label{lem:distinct-pairs}
If \(u\ne v\), then the defining pairs of \(u\) and \(v\) are distinct.  Moreover,
if \(c(u)=c(v)\), then \(\Bags(u)\cap\Bags(v)=\emptyset\); in particular, \(u\)
and \(v\) are non-adjacent.
\end{lemma}

\begin{proof}
If two distinct vertices had the same defining pair \((U,i)\), choose any
\(t\in U\).  Both vertices would lie in \(\beta(t)\) and have color \(i\),
contradicting bag-injectivity.  Thus their defining pairs are distinct.

Now assume \(c(u)=c(v)=i\).  If \(t\in\Bags(u)\cap\Bags(v)\), then both vertices
lie in \(\beta(t)\) with color \(i\), again contradicting bag-injectivity.  Hence
\(\Bags(u)\cap\Bags(v)=\emptyset\).  If \(uv\) were an edge, the edge-covering
axiom of the tree-decomposition would give a bag containing both \(u\) and \(v\),
which is impossible.  Therefore \(u\) and \(v\) are non-adjacent.
\end{proof}

\begin{lemma}[Edges from common bags]\label{lem:edges-bags}
For vertices \(u,v\in V(G)\), the edge \(uv\) belongs to \(G\) iff there exists a
node \(t\in\Bags(u)\cap\Bags(v)\) such that the colors \(c(u),c(v)\) are adjacent
in the node graph \(H_t\).
\end{lemma}

\begin{proof}
If \(uv\in E(G)\), then some bag \(\beta(t)\) contains both endpoints.  Hence
\(t\in\Bags(u)\cap\Bags(v)\), and the induced node graph \(H_t=G[\beta(t)]\)
contains the edge between the two color representatives \(c(u)\) and \(c(v)\).

Conversely, if such a node \(t\) exists, the letter \(H_t\) records an edge
between the unique vertices of colors \(c(u)\) and \(c(v)\) in \(\beta(t)\).  Since
\(H_t\) is an induced subgraph of \(G\), that edge is an edge of \(G\).  Those
unique vertices are exactly \(u\) and \(v\), because \(t\in\Bags(u)\cap\Bags(v)\).
\end{proof}

A vertex is represented on a \(\Sigma\)-tree by a tuple of node sets rather than
by a color parameter.  The \emph{defining tuple} of a vertex \(v\) is
\(\tuple U^v=(U^v_0,\ldots,U^v_\omega)\), where
\[
  U^v_i =
  \begin{cases}
    \Bags(v), & i=c(v),\\
    \emptyset, & i\ne c(v).
  \end{cases}
\]
For a vertex set \(S\subseteq V(G)\), its defining tuple \(\tuple U^S\) is given by
\[
  U^S_i = \bigcup_{v\in S,\ c(v)=i}\Bags(v).
\]
The unions are disjoint by Lemma~\ref{lem:distinct-pairs}.

\begin{lemma}[Defining tuples of vertex sets]\label{lem:defining-tuples}
A tuple \((U_0,\ldots,U_\omega)\) is the defining tuple of a vertex set \(S\subseteq V(G)\) iff,
for each color \(i\), the set \(U_i\) can be partitioned into nonempty pieces
\((U_{i,a})_{a\in A_i}\) such that each pair \((U_{i,a},i)\) satisfies the five
conditions of Lemma~\ref{lem:defining-pairs}.
\end{lemma}

\begin{proof}
If \(\tuple U\) is the defining tuple of \(S\), then for each color \(i\), the
pieces are \(\Bags(v)\) for vertices \(v\in S\) with color \(i\).  Lemma~\ref{lem:defining-pairs} shows
that every piece satisfies the five conditions, and Lemma~\ref{lem:distinct-pairs} shows that pieces of
the same color are disjoint.

Conversely, suppose such partitions are given.  By Lemma~\ref{lem:defining-pairs}, each piece
\((U_{i,a},i)\) is the defining pair of a unique vertex \(v_{i,a}\) of color
\(i\).  Let
\[
  S\defeq\{v_{i,a}:0\le i\le\omega,\ a\in A_i\}.
\]
For each \(i\), the union of \(\Bags(v_{i,a})\) over \(a\in A_i\) is exactly the partition union \(U_i\).  Therefore
\(\tuple U\) is the defining tuple of \(S\).
\end{proof}

\section{Converting graph \texorpdfstring{\(\MSO\)}{MSO} to tree \texorpdfstring{\(\MSO\)}{MSO}}

All formulas in this section are over
\(\tau_\Sigma=\{\child_1,\child_2\}\cup\{P_a:a\in\Sigma_\omega\}\).  Since
\(\Sigma_\omega\) is finite, every disjunction over a subset of \(\Sigma_\omega\)
is a finite \MSO{} formula.

\subsection{Basic tree formulas}

Use the following standard abbreviations:
\[
  \parent(y,x) \defeq \child_1(y,x)\vee\child_2(y,x),
  \qquad
  \rootnode(x) \defeq \neg\exists y\,\parent(y,x).
\]
Set inclusion, emptiness, and nonemptiness are expressed by
\[
  \subsetf(X,Y)\defeq \forall x(x\in X\to x\in Y),\quad
  \emptysetf(X)\defeq \forall x\,x\notin X,
\]
\[
  \nonempty(X)\defeq \exists x\,x\in X.
\]
Let \(\operatorname{Adj}_T(x,y)\) abbreviate
\(\parent(x,y)\vee\parent(y,x)\).  A set \(X\) is connected if it is nonempty and
no nontrivial partition of \(X\) has no tree edge across it:
\[
\begin{aligned}
  \conn(X) \defeq{}& \nonempty(X)\wedge \\
  &\forall Y\Bigl(
    \subsetf(Y,X)\wedge\nonempty(Y)\wedge\exists z(z\in X\wedge z\notin Y)
    \to \\
  &\qquad\qquad \exists u\exists v\,
    (u\in Y\wedge v\in X\wedge v\notin Y\wedge \operatorname{Adj}_T(u,v))
  \Bigr).
\end{aligned}
\]
For connected \(X\), the top node is the unique node of \(X\) whose parent is not
in \(X\):
\[
\begin{aligned}
  \topnode(x,X) \defeq{}& x\in X\wedge \conn(X)\wedge
  \bigl(\rootnode(x)\vee \exists p(\parent(p,x)\wedge p\notin X)\bigr),\\
  \dang(x,X) \defeq{}& x\notin X\wedge \exists p(\parent(p,x)\wedge p\in X).
\end{aligned}
\]
In a rooted tree, a nonempty connected set has exactly one node satisfying
\(\topnode(x,X)\).

For a color \(i\in C_\omega\), define finite letter sets
\[
  \Sigma_i\defeq\{(H,R):i\in V(H)\},\qquad
  \Sigma_i^R\defeq\{(H,R):i\in R\}.
\]
For colors \(i,j\), define
\[
  \Sigma_{ij}^E\defeq\{(H,R): ij\in E(H)\}.
\]
Let
\[
  A_i(x)\defeq \bigvee_{a\in\Sigma_i} P_a(x),\quad
  R_i(x)\defeq \bigvee_{a\in\Sigma_i^R} P_a(x),\quad
  E_{ij}(x)\defeq \bigvee_{a\in\Sigma_{ij}^E} P_a(x).
\]
For \(Q\in\Pred\), define
\[
  \Sigma_i^Q\defeq\{(H,R): i\in Q^H\},
  \qquad
  Q_i(x)\defeq\bigvee_{a\in\Sigma_i^Q}P_a(x).
\]

\subsection{Recognizing vertices and vertex sets}

Define
\[
\begin{aligned}
  \varphi_{\vtx,i}(Z) \defeq{}&
  \varphi_{\legal}\wedge \conn(Z)\\
  &\wedge\ \forall x(x\in Z\to A_i(x))\\
  &\wedge\ \forall x((x\in Z\wedge\neg\topnode(x,Z))\to R_i(x))\\
  &\wedge\ \forall x(\topnode(x,Z)\to \neg R_i(x))\\
  &\wedge\ \forall x(\dang(x,Z)\to \neg R_i(x)).
\end{aligned}
\]

\begin{lemma}\label{lem:vtx-i}
For every \(i\in C_\omega\), the formula \(\varphi_{\vtx,i}(Z)\) holds in a
\(\Sigma\)-tree \((T,\rho)\) with parameter \(U\subseteq N(T)\) iff \((T,\rho)\) is
legal and \((U,i)\) is the defining pair of a vertex of the decoded graph.
\end{lemma}

\begin{proof}
The conjunct \(\varphi_{\legal}\) gives legality.  The remaining conjuncts are a
literal \MSO{} translation of the five conditions in Lemma~\ref{lem:defining-pairs}: connectedness of
\(U\); existence of color \(i\) in every bag of \(U\); existence of color \(i\) in
each non-top adhesion; absence of color \(i\) in the top adhesion; and absence of
color \(i\) in every dangling child adhesion.  Therefore, by Lemma~\ref{lem:defining-pairs}, the
formula holds exactly for defining pairs.
\end{proof}

For a tuple \(\tuple Z=(Z_0,\ldots,Z_\omega)\), define
\[
  \varphi_{\vtx}(\tuple Z)
  \defeq
  \bigvee_{i=0}^{\omega}
  \left(
    \varphi_{\vtx,i}(Z_i)\wedge
    \bigwedge_{j\ne i}\emptysetf(Z_j)
  \right).
\]

\begin{lemma}
There is an \MSO{} formula \(\varphi_{\vtx}(\tuple Z)\) with \(\omega+1\) free set
variables such that \((T,\rho)\models\varphi_{\vtx}(\tuple U)\) iff \((T,\rho)\) is legal and \(\tuple U\) is
the defining tuple of a vertex in the decoded graph.
\end{lemma}

\begin{proof}
If \(\tuple U\) is the defining tuple of a vertex \(v\) of color \(i\), then
\(U_i=\Bags(v)\) and all other coordinates are empty, so the \(i\)-th disjunct
holds by Lemma~\ref{lem:vtx-i}.  Conversely, if the formula holds, some coordinate \(U_i\) is
a defining pair of a color-\(i\) vertex and all other coordinates are empty;
therefore \(\tuple U\) is exactly that vertex's defining tuple.
\end{proof}

Define
\[
  \varphi_{\set}(\tuple Z)
  \defeq
  \varphi_{\legal}\wedge
  \bigwedge_{i=0}^{\omega}
  \forall x\left(
    x\in Z_i\to
    \exists Y\bigl(\subsetf(Y,Z_i)\wedge x\in Y\wedge\varphi_{\vtx,i}(Y)\bigr)
  \right).
\]

\begin{lemma}
There is an \MSO{} formula \(\varphi_{\set}(\tuple Z)\) with \(\omega+1\) free set
variables such that \((T,\rho)\models\varphi_{\set}(\tuple U)\) iff \((T,\rho)\) is
legal and \(\tuple U\) is the defining tuple of a vertex subset of the decoded
graph.
\end{lemma}

\begin{proof}
Assume \(\tuple U\) is the defining tuple of a set \(S\).  If
\(x\in U_i\), then \(x\in\Bags(v)\) for a unique vertex \(v\in S\) of color \(i\);
using \(Y=\Bags(v)\) satisfies the displayed formula by Lemma~\ref{lem:vtx-i}.

Conversely, suppose \(\varphi_{\set}(\tuple U)\) holds.  For every node
\(x\in U_i\), choose a set \(Y_x\subseteq U_i\) containing \(x\) with
\(\varphi_{\vtx,i}(Y_x)\).  By Lemma~\ref{lem:vtx-i}, each \((Y_x,i)\) is the defining pair of
a color-\(i\) vertex.  Distinct color-\(i\) vertices have disjoint bags by
Lemma~\ref{lem:distinct-pairs}; hence the distinct sets among the \(Y_x\)'s form a partition of
\(U_i\) into pieces satisfying Lemma~\ref{lem:defining-pairs}.  Lemma~\ref{lem:defining-tuples} then gives a vertex set whose
defining tuple is \(\tuple U\).
\end{proof}

\subsection{Recognizing graph and predicate atomic formulas}

For two tuples \(\tuple X=(X_0,\ldots,X_\omega)\) and \(\tuple Y=(Y_0,\ldots,Y_\omega)\), define
\[
\begin{aligned}
  \varphi_{\adj}(\tuple X,\tuple Y)\defeq{}&
  \varphi_{\vtx}(\tuple X)\wedge\varphi_{\vtx}(\tuple Y)\\
  &\wedge\bigvee_{i,j=0}^{\omega}
  \exists t\bigl(t\in X_i\wedge t\in Y_j\wedge E_{ij}(t)\bigr).
\end{aligned}
\]

\begin{lemma}
The formula \(\varphi_{\adj}(\tuple X,\tuple Y)\) holds on defining tuples
\(\tuple U^u,\tuple U^v\) iff \(uv\) is an edge in the decoded graph.
\end{lemma}

\begin{proof}
The first two conjuncts ensure that \(\tuple X\) and \(\tuple Y\) define vertices,
say \(u\) and \(v\).  If the last conjunct holds for colors \(i=c(u)\) and
\(j=c(v)\), then there is a node \(t\in\Bags(u)\cap\Bags(v)\) whose letter records
an edge between colors \(i,j\).  Lemma~\ref{lem:edges-bags} implies \(uv\in E(G)\).  Conversely, if
\(uv\in E(G)\), Lemma~\ref{lem:edges-bags} gives such a node \(t\), so the formula holds.
\end{proof}

For \(Q\in\Pred\), define
\[
  \varphi_Q(\tuple X)\defeq
  \varphi_{\vtx}(\tuple X)\wedge
  \bigvee_{i=0}^{\omega}\exists t\bigl(t\in X_i\wedge Q_i(t)\bigr).
\]

\begin{lemma}
Let \(Q\in\Pred\).  The formula \(\varphi_Q(\tuple X)\) holds on the defining
tuple \(\tuple U^v\) iff \(v\in Q^G\).
\end{lemma}

\begin{proof}
If \(\tuple X=\tuple U^v\), then exactly one coordinate is nonempty, namely
\(\Bags(v)\) at the color \(i=c(v)\).  In an encoding, every bag containing \(v\)
records the same unary predicate information for the color \(i\), because each
node letter is induced from the same ambient \(\tau_{\Pred}\)-graph \(G\).
Therefore some, equivalently every, node \(t\in\Bags(v)\) satisfies \(Q_i(t)\)
exactly when \(v\in Q^G\).
\end{proof}

Define set equality by
\[
  X=Y \quad\text{as an abbreviation for}\quad
  \forall z(z\in X\leftrightarrow z\in Y),
\]
and put
\[
  \varphi_{\eq}(\tuple X,\tuple Y)
  \defeq
  \varphi_{\vtx}(\tuple X)\wedge\varphi_{\vtx}(\tuple Y)
  \wedge\bigwedge_{i=0}^{\omega} X_i=Y_i.
\]

\begin{lemma}
The formula \(\varphi_{\eq}(\tuple X,\tuple Y)\) holds on defining tuples
\(\tuple U^u,\tuple U^v\) iff \(u=v\) in the decoded graph.
\end{lemma}

\begin{proof}
If \(u=v\), the defining tuples are identical.  Conversely, if the tuples are
coordinatewise equal, then the nonempty coordinate and the corresponding bag set
are the same for both vertices.  Thus the defining pairs are equal.  By
Lemma~\ref{lem:distinct-pairs}, distinct vertices have distinct defining pairs, so \(u=v\).
\end{proof}

For containment, define
\[
\begin{aligned}
  \varphi_{\cont}(\tuple X,\tuple Y)\defeq{}&
  \varphi_{\vtx}(\tuple X)\wedge\varphi_{\set}(\tuple Y)\\
  &\wedge\bigvee_{i=0}^{\omega}
  \bigl(\nonempty(X_i)\wedge\subsetf(X_i,Y_i)\bigr).
\end{aligned}
\]

\begin{lemma}
The formula \(\varphi_{\cont}(\tuple X,\tuple Y)\) holds iff \(\tuple X\) is the
defining tuple of a vertex \(v\), \(\tuple Y\) is the defining tuple of a vertex
set \(S\), and \(v\in S\).
\end{lemma}

\begin{proof}
Let \(i=c(v)\).  If \(v\in S\), then \(X_i=\Bags(v)\) is one of the pieces whose
union forms \(Y_i\), so \(X_i\subseteq Y_i\), and the formula holds.

Conversely, suppose the formula holds.  Then for the unique nonempty coordinate
\(X_i=\Bags(v)\), we have \(X_i\subseteq Y_i\).  Since \(\tuple Y\) defines a set
\(S\), the coordinate \(Y_i\) is the disjoint union of \(\Bags(w)\) over
\(w\in S\) of color \(i\).  The connected set \(\Bags(v)\) is contained in this
disjoint union and is itself a defining pair.  By Lemma~\ref{lem:distinct-pairs}, it cannot intersect
any different color-\(i\) defining pair.  Hence one of the pieces is
\(\Bags(v)\), so \(v\in S\).
\end{proof}

\subsection{The formula translation}

Let \(\theta\) be an \MSO{} formula over the unary-expanded graph vocabulary
\(\tau_{\Pred}\).  Each first-order variable \(x\) is represented by a tuple of
set variables
\[
  \tuple Z_x=(Z^x_0,\ldots,Z^x_\omega),
\]
and each monadic set variable \(X\) is represented by
\[
  \tuple Z_X=(Z^X_0,\ldots,Z^X_\omega).
\]
Define the translation \(\theta^\star\) recursively:
\begin{align*}
  (x=y)^\star &\defeq \varphi_{\eq}(\tuple Z_x,\tuple Z_y),\\
  (\adj(x,y))^\star &\defeq \varphi_{\adj}(\tuple Z_x,\tuple Z_y),\\
  (Q(x))^\star &\defeq \varphi_Q(\tuple Z_x)\qquad(Q\in\Pred),\\
  (x\in X)^\star &\defeq \varphi_{\cont}(\tuple Z_x,\tuple Z_X),
\end{align*}
and commute the translation with Boolean connectives.  Existential quantifiers
are translated by
\[
  (\exists x\,\psi)^\star
  \defeq
  \exists Z^x_0\cdots\exists Z^x_\omega
  \bigl(\varphi_{\vtx}(\tuple Z_x)\wedge \psi^\star\bigr),
\]
\[
  (\exists X\,\psi)^\star
  \defeq
  \exists Z^X_0\cdots\exists Z^X_\omega
  \bigl(\varphi_{\set}(\tuple Z_X)\wedge \psi^\star\bigr).
\]
Universal quantifiers are handled by the identity \(\forall q\,\psi\equiv
\neg\exists q\,\neg\psi\).  For a sentence \(\varphi\), define
\[
  \varphi_{\mathrm{tree}}\defeq \varphi_{\legal}\wedge\varphi^\star.
\]

\paragraph{Correctness of the translation.}
Let \((T,\rho)\) be the encoding of \((G,c,(T,\beta))\).  If a graph assignment sends every
individual variable to a vertex and every set variable to a vertex set, let the
corresponding tree assignment send each variable to the defining tuple of that
object.  A structural induction on \(\theta\), using the atomic correctness
lemmas above for equality, adjacency, unary predicates, and membership, together
with the displayed clauses for quantifiers, gives
\[
  G\models\theta
  \quad\Longleftrightarrow\quad
  (T,\rho)\models\theta^\star
\]
under corresponding assignments.  In particular, for every \(\tau_{\Pred}\)-MSO{}
sentence \(\varphi\),
\[
  G\models\varphi
  \quad\Longleftrightarrow\quad
  (T,\rho)\models\varphi_{\mathrm{tree}}.
\]

\section{Completing the proof of Courcelle's theorem}

\begin{proof}[Proof of Theorem~\ref{thm:courcelle}]
Fix \(\omega\) and the finite unary vocabulary \(\Pred\).  Let \(G\) be a
\(\tau_{\Pred}\)-graph, let \((T,\beta)\) be a binary tree-decomposition of its
underlying graph of width at most \(\omega\), and let \(\varphi\) be an \MSO{}
sentence over \(\tau_{\Pred}\).
Choose a bag-injective coloring \(c:V(G)\to C_\omega\).  Construct the encoding
\((T,\rho)\) over the finite alphabet \(\Sigma_\omega\); this is linear in the
number of nodes of the decomposition because every bag has bounded size and
\(\Pred\) is fixed.

Translate \(\varphi\) to the tree sentence \(\varphi_{\mathrm{tree}}
=\varphi_{\legal}\wedge\varphi^\star\).  The size of this sentence depends only
on \(\omega\), \(\Pred\), and \(\varphi\), because all alphabet disjunctions are
over the finite set \(\Sigma_\omega\).  By the correctness of the translation,
\[
  G\models\varphi
  \quad\Longleftrightarrow\quad
  (T,\rho)\models\varphi_{\mathrm{tree}}.
\]

Use the standard theorem that \MSO{}-definable languages of finite labelled
binary trees are exactly regular tree languages.  Thus there is a deterministic
bottom-up finite tree automaton \(\Aut\), depending only on \(\omega\), \(\Pred\),
and \(\varphi\), which accepts exactly the \(\Sigma_\omega\)-trees satisfying
\(\varphi_{\mathrm{tree}}\).  Running \(\Aut\) on \((T,\rho)\) is linear in
\(|N(T)|\).  Accept iff \(\Aut\) accepts.  The displayed equivalence proves
correctness, and the automaton size and translation cost are absorbed into
\(f_{\Pred}(\omega,|\varphi|)\).  After the usual linear-size normalization of
the decomposition, this gives the advertised linear dependence on \(|V(G)|\)
for fixed \(\varphi\).
\end{proof}

\begin{remark}
The only external ingredient in the last paragraph is the regularity theorem for
\MSO{} on finite labelled binary trees.  All graph-specific work is contained in
the encoding, decoding, and interpretation lemmas proved above.
\end{remark}

\begin{thebibliography}{9}

\bibitem{Courcelle1990}
Bruno Courcelle.
\newblock The monadic second-order logic of graphs. I. Recognizable sets of
finite graphs.
\newblock \emph{Information and Computation}, 85(1):12--75, 1990.

\bibitem{Korhonen2022}
Tuukka Korhonen.
\newblock A single-exponential time 2-approximation algorithm for treewidth.
\newblock In \emph{2021 IEEE 62nd Annual Symposium on Foundations of Computer
Science (FOCS)}, pages 184--192, 2022.

\bibitem{Bojanczyk2015}
Mikolaj Bojanczyk.
\newblock Interactive lecture note: Interpreting a graph in its tree decomposition.
\newblock 2015.

\end{thebibliography}

\end{document}
